\SourceLeaf{322}
\[
\begin{array}{rcll@{\qquad}rcll}
2 & \text{um} & 180^0+A & \text{rechts von }1 & 3 & \text{um} & B & \text{rechts von }1\\[2pt]
 & \text{,,} & A & \text{,,} & & \text{,,} & 180^0+B & \text{,,}\\[2pt]
 & \text{,,} & 180^0+A & \text{,,} & & \text{,,} & 180^0+B & \text{,,}\;.
\end{array}
\]

Das Resultat selbst aber quadrirt immer nur mit Einem dieser vier Systeme von Datis, und die andern drei sind \textls[120]{physisch unmöglich}.

Der metaphysische Grund dieser Erscheinung ist, dass man von der beobachteten Richtung weiter nichts benutzt, als dass die \textls[120]{geraden Linien} gewisse Winkel mit einander machen, indem man diese geraden Linien auf \textls[120]{beiden} Seiten in indefinitum sich erstrecken lässt, während das Fortschreiten des Lichts nur in Einem bestimmten Sinne geschieht, also die Fälle, wo der Gegenstand rückwärts läge, ausgeschlossen werden müssen.

Zur Vollständigkeit gehört also die Auflösung der Aufgabe, den Datis vor geführter Rechnung gleich anzusehen, ob sie vereinbar sind. Es lässt sich diess sehr zierlich abthun, ich will aber Ihnen nicht vorgreifen, es Sich selbst zu entwickeln. \mbox{. . . . . . .}



\centerline{\textsc{Gauss} an \textsc{Gerling}. Göttingen, 14. Januar 1842.}

\indent \mbox{. . . . . . . .} Das Criterium wegen der \textsc{Pothenot}schen Aufgabe theile ich Ihnen zwar gern mit, wünsche jedoch, dass Sie es für Sich behalten, und zwar aus dem Grunde, weil ich das Theorem, womit es zusammenhängt, selbst einmal bei schicklicher Gelegenheit zu behandeln mir vorbehalte, weniger wegen der Eleganz des Theorems an sich, als wegen der Eleganz, welche die Anwendung der complexen Grössen dabei darbietet; also namentlich bei einer Gelegenheit, wo ich mehr von dem Gebrauch der complexen Grössen sagen kann. Würde das Theorem vorher anderweit ins Publicum gebracht, so würde für mich nachher keine Lust mehr bleiben, selbst darauf zurückzukommen. Also jetzt zur Sache.

Man kann das Criterium auf vielfache Weise einkleiden; am zierlichsten wird es sein, es so zu thun, dass alle drei gegebenen Punkte symmetrisch dabei vorkommen. Ich bezeichne also die drei gegebenen Punkte mit 1, 2, 3;
